%==============================================================================
%  自定义命令与数学符号
%  —— 统一命名，避免同一个符号在不同章节写法不一致
%==============================================================================

% --- 常用量（用 \ensuremath 包裹，正文与公式里都能直接写）---
\newcommand{\Np}{\ensuremath{N}}                   % 无人机数量
\newcommand{\Mt}{\ensuremath{M}}                   % 目标数量
\newcommand{\CR}{\ensuremath{\mathrm{CR}}}         % 交叉概率
\newcommand{\Ffactor}{\ensuremath{F}}              % 缩放因子
\newcommand{\NP}{\ensuremath{\mathrm{NP}}}         % 种群规模
\newcommand{\Gmax}{\ensuremath{G_{\max}}}          % 最大代数
\newcommand{\rhoinit}{\ensuremath{\rho}}           % PopInit 注入比例
\newcommand{\vx}{\ensuremath{\bm{x}}}              % 决策向量
\newcommand{\vtheta}{\ensuremath{\bm{\theta}}}     % 参数向量
\newcommand{\Expect}[1]{\ensuremath{\mathbb{E}\left[#1\right]}}

% --- 方法名（保持全文大小写一致）---
\newcommand{\dmde}{\textsc{dmde}}
\newcommand{\llmdmde}{\textsc{llm-dmde}}
\newcommand{\popinit}{\textsc{PopInit}}
\newcommand{\crctrl}{\textsc{CR-Control}}

% --- 审阅辅助：编译后显示为红色，投稿前用下面的开关一键隐藏 ---
\newcommand{\reviewmode}{1}           % 投稿前改成 0 即可全部隐藏
\ifnum\reviewmode=1
  % 注意：这里刻意不加 \textbf / \itshape——\popinit 等命令用 \textsc，
  % 而 cm 字体在粗体下没有 small-caps 字型，会触发 "Font shape undefined" 警告
  \newcommand{\todo}[1]{\textcolor{red}{[TODO: #1]}}
  \newcommand{\note}[1]{\textcolor{blue}{\textit{[#1]}}}
\else
  \newcommand{\todo}[1]{}
  \newcommand{\note}[1]{}
\fi

% --- 图表引用（Elsevier 惯例：Fig. / Table，无论句首句中都这样写）---
\newcommand{\figref}[1]{Fig.~\ref{#1}}
\newcommand{\tabref}[1]{Table~\ref{#1}}
\newcommand{\equref}[1]{Eq.~\eqref{#1}}
\newcommand{\secref}[1]{Section~\ref{#1}}
